% ===============================================================
% CHAPITRE 3 : SPÉCIFICATION DES BESOINS ET CONCEPTION DU SYSTÈME
% ===============================================================

\chapter{Spécification des Besoins et Conception du Système}

\section*{Introduction}

Ce chapitre formalise les besoins du système WC Fantasy, identifie les acteurs, définit le cahier des charges et présente la conception complète : architecture globale, diagrammes UML (cas d'utilisation, activité), schéma de base de données et choix technologiques justifiés.

\section{Analyse des besoins}

\subsection{Besoins fonctionnels}

\vspace{0.3cm}

\begin{table}[!htbp]
\centering
    \begin{tabularx}{\linewidth}{|l|X|}
        \hline
        \rowcolor{emsigreen}\textcolor{white}{\textbf{\#}} & \textcolor{white}{\textbf{Besoin fonctionnel}} \\
        \hline
        BF01 & S'inscrire, se connecter et gérer son profil utilisateur \\
        \hline
        BF02 & Consulter le catalogue complet des joueurs de la CdM 2026 avec stats et photos \\
        \hline
        BF03 & Composer une équipe de 15 joueurs dans un budget de 100M avec contraintes de positions \\
        \hline
        BF04 & Créer ou rejoindre une ligue privée ou publique via code d'invitation \\
        \hline
        BF05 & Effectuer des transferts de joueurs (entrant/sortant) en période de transfert \\
        \hline
        BF06 & Participer à des batailles hebdomadaires head-to-head contre d'autres équipes de la ligue \\
        \hline
        BF07 & Consulter les scores et classements de la ligue en temps réel \\
        \hline
        BF08 & Soumettre des prédictions communautaires avant chaque journée (top buteur, buts, nation) \\
        \hline
        BF09 & Communiquer avec les membres de la ligue via un chat temps réel \\
        \hline
        BF10 & Interagir avec un chatbot IA pour obtenir des informations et conseils \\
        \hline
        BF11 & Gérer les actualités football (admin) et les consulter (utilisateur) \\
        \hline
        BF12 & Compléter des défis quotidiens/hebdomadaires pour gagner des points XP \\
        \hline
        BF13 & Réinitialiser son mot de passe par email (lien sécurisé, TTL 1h) \\
        \hline
        BF14 & Administrer : joueurs, scores de GW, utilisateurs, prédictions (rôle admin) \\
        \hline
    \end{tabularx}
\caption{Besoins fonctionnels du système}
\end{table}

\vspace{0.2cm}
\noindent\small\textit{Ce tableau énumère les besoins fonctionnels (BF) requis pour le bon fonctionnement de l'application, couvrant les fonctionnalités utilisateurs et administratives.}\normalsize

\subsection{Besoins non fonctionnels}

\vspace{0.3cm}

\begin{table}[!htbp]
\centering
    \begin{tabularx}{\linewidth}{|l|X|}
        \hline
        \rowcolor{emsigreen}\textcolor{white}{\textbf{Catégorie}} & \textcolor{white}{\textbf{Exigence}} \\
        \hline
        Performance & Temps de réponse API $< 200$ ms pour 95\% des requêtes grâce au cache Redis \\
        \hline
        Scalabilité & Architecture stateless permettant la multiplication des instances backend \\
        \hline
        Sécurité & JWT avec rotation de refresh token, hachage bcrypt (10 rounds), rate limiting \\
        \hline
        Disponibilité & Reconnexion automatique Socket.io, retry interceptor côté Flutter \\
        \hline
        Fiabilité & Transactions PostgreSQL ACID pour toutes les opérations financières (budget) \\
        \hline
        Maintenabilité & Architecture MVC côté backend, code modulaire et centralisé (pointsCalculator) \\
        \hline
        Compatibilité & Application Flutter compatible Android 5.0+ et iOS 12+ \\
        \hline
        Accessibilité & Interface intuitive, temps d'apprentissage $<$ 5 minutes pour un nouvel utilisateur \\
        \hline
    \end{tabularx}
\caption{Exigences non fonctionnelles du système}
\end{table}

\vspace{0.2cm}
\noindent\small\textit{Ce tableau récapitule les exigences non fonctionnelles de sécurité, de performance, de scalabilité et de compatibilité nécessaires pour garantir la qualité de service de la plateforme.}\normalsize

\section{Identification des acteurs}

Le système WC Fantasy implique deux acteurs principaux :

\begin{itemize}
    \item \textbf{Utilisateur (User)} : tout membre inscrit sur la plateforme. Il peut composer son équipe, rejoindre/créer des ligues, effectuer des transferts, participer aux batailles, prédire, chatter et interagir avec le chatbot IA.
    \item \textbf{Administrateur (Admin)} : utilisateur disposant de droits étendus. Il gère les joueurs du catalogue, saisit les scores réels après chaque journée de jeu, publie des actualités, résout les prédictions communautaires et administre les utilisateurs.
\end{itemize}

\section{Cahier des charges synthétique}

\vspace{0.3cm}

\begin{table}[!htbp]
\centering
    \begin{tabularx}{\linewidth}{|l|X|}
        \hline
        \rowcolor{emsigreen}\multicolumn{2}{|c|}{\textcolor{white}{\textbf{Récapitulatif des choix technologiques}}} \\
        \hline
        Backend & Node.js + Express.js --- API REST\cite{mdn_rest} --- PostgreSQL --- Redis --- Socket.io \\
        \hline
        Mobile & Flutter (Dart) --- Riverpod (gestion état) --- Dio (HTTP) --- FlutterSecureStorage \\
        \hline
        Auth & JWT (access 7j + refresh 30j) --- bcrypt --- rate limiting --- CORS configurable \\
        \hline
        Temps réel & Socket.io --- rooms par ligue et par bataille --- rate limit 1 msg/sec \\
        \hline
        Scoring & Barème détaillé par position --- synchronisation auto via football-data.org \\
        \hline
        Gamification & XP, défis quotidiens/hebdomadaires, achievements débloquables \\
        \hline
        IA & Chatbot Google Gemini --- module prédictions communautaires \\
        \hline
        Email & Brevo SMTP --- reset mot de passe --- email de bienvenue \\
        \hline
    \end{tabularx}
\caption{Récapitulatif des choix technologiques}
\end{table}

\vspace{0.2cm}
\noindent\small\textit{Ce tableau récapitule les choix technologiques de l'architecture globale en listant les outils utilisés pour le backend, l'application mobile, le temps réel et l'administration.}\normalsize

\section{Conception générale --- Architecture globale}

L'architecture du système repose sur un modèle client-serveur découplé en trois couches principales. Le schéma ci-dessous illustre l'organisation générale du système et les interactions entre ses composants :

\vspace{0.5cm}

\begin{figure}[H]
    \centering
    \includegraphics[width=0.85\textwidth]{photos/global_architecture.jpg}    \caption{Architecture globale du système WC Fantasy}
    \label{fig:architecture}
\end{figure}

\vspace{0.2cm}
\noindent\small\textit{Cette figure schématise l'interaction entre l'application mobile (Flutter), le serveur backend (Node.js/Express) et les différentes bases de données. Elle montre comment les requêtes HTTP et les événements WebSocket transitent pour fournir une expérience temps réel à l'utilisateur.}
\normalsize

\vspace{0.5cm}

L'architecture se décompose en trois couches distinctes :

\begin{itemize}
    \item \textbf{Couche Cliente (Flutter)} : interface utilisateur cross-platform, gestion de l'état avec Riverpod, communication HTTP via Dio avec intercepteurs d'authentification et retry automatique.
    \item \textbf{Couche Serveur (Node.js/Express)} : logique métier organisée en Controllers, Routes et Middlewares, serveur Socket.io intégré, tâche cron de synchronisation des scores, service email.
    \item \textbf{Couche Données} : PostgreSQL pour la persistance (20 tables), Redis pour le cache (TTL configurable), APIs externes (football-data.org, Google Gemini, Brevo SMTP, Wikipedia).
\end{itemize}

\section{Modélisation fonctionnelle --- Diagramme de cas d'utilisation}

Le diagramme de cas d'utilisation ci-dessous recense l'ensemble des interactions entre les deux acteurs (Utilisateur et Administrateur) et le système. Il est organisé par domaine fonctionnel.

\vspace{0.5cm}

\begin{figure}[H]
    \centering
    \includegraphics[width=0.6\textwidth]{photos/Diagramme_utilisateur.png}
    \caption{Diagramme de cas d'utilisation --- Utilisateur}
    \label{fig:usecase_utilisateur}
\end{figure}

\vspace{0.2cm}
\noindent\small\textit{Ce diagramme affiche toutes les actions réalisables par un utilisateur standard, telles que l'inscription, la gestion de l'équipe, la participation aux batailles et l'interaction avec le chatbot IA.}
\normalsize

\vspace{0.5cm}

\begin{figure}[H]
    \centering
    \includegraphics[width=0.6\textwidth]{photos/Diagramme_admin.png}
    \caption{Diagramme de cas d'utilisation --- Administrateur}
    \label{fig:usecase_admin}
\end{figure}

\vspace{0.2cm}
\noindent\small\textit{Ce diagramme présente les privilèges de l'administrateur, incluant la gestion du catalogue des joueurs, la validation des scores après chaque journée et la résolution des prédictions communautaires.}
\normalsize

\vspace{0.5cm}

Le diagramme met en évidence cinq grands domaines fonctionnels :

\begin{itemize}
    \item \textbf{Authentification} : inscription, connexion, réinitialisation du mot de passe.
    \item \textbf{Équipe} : composer son équipe, effectuer des transferts, désigner un capitaine.
    \item \textbf{Compétition} : créer/rejoindre des ligues, participer aux batailles, consulter les journées.
    \item \textbf{Profil \& Joueurs} : consulter le profil d'un joueur, suivre ses statistiques.
    \item \textbf{Information \& Assistance} : actualités, chatbot IA.
    \item \textbf{Administration} : gérer les utilisateurs, joueurs, scores, ligues et prédictions.
\end{itemize}

\section{Modélisation structurelle --- Schéma de base de données}

La base de données PostgreSQL du projet comporte \textbf{20 tables} organisées autour de l'entité centrale \texttt{users}. Le tableau ci-dessous présente les tables principales avec leur rôle et les colonnes clés :

\vspace{0.3cm}

\begin{table}[!htbp]
\centering
    \begin{tabularx}{\linewidth}{|l|X|}
        \hline
        \rowcolor{emsigreen}\textcolor{white}{\textbf{Table}} & \textcolor{white}{\textbf{Rôle \& Colonnes principales}} \\
        \hline
        users & Comptes utilisateurs : id, username, email, password\_hash, total\_points, xp, rank \\
        \hline
        teams & Équipes fantasy : id, user\_id, league\_id, budget, formation, captain\_id, players (JSONB) \\
        \hline
        players & Catalogue joueurs CdM : id, name, position, club, nationality, price, points, stats (JSONB) \\
        \hline
        leagues & Ligues : id, name, owner\_id, max\_members, is\_private, invite\_code, season, status \\
        \hline
        league\_members & Appartenance ligue-utilisateur : league\_id, user\_id, points \\
        \hline
        battles & Batailles H2H : id, league\_id, home\_team\_id, away\_team\_id, home\_score, away\_score, status \\
        \hline
        gameweeks & Journées de jeu : id, number, season, status, deadline \\
        \hline
        player\_scores & Scores réels par journée : goals, assists, clean\_sheet, saves, yellow\_cards, total\_points \\
        \hline
        predictions & Prédictions : predicted\_top\_scorer, predicted\_winner\_nation, predicted\_goals, bonus\_points \\
        \hline
        chat\_messages & Messages du chat de ligue : league\_id, user\_id, message, created\_at \\
        \hline
        refresh\_tokens & Tokens de rafraîchissement JWT : user\_id, token, expires\_at \\
        \hline
        user\_challenges & Défis complétés : user\_id, challenge\_key, challenge\_type, xp\_earned \\
        \hline
        user\_achievements & Achievements débloqués : user\_id, achievement\_key, unlocked\_at \\
        \hline
        password\_reset\_tokens & Tokens de reset MDP : user\_id, token\_hash, expires\_at \\
        \hline
    \end{tabularx}
\caption{Schéma de la base de données}
\end{table}

\vspace{0.2cm}
\noindent\small\textit{Ce tableau décrit l'architecture relationnelle de la base de données PostgreSQL en listant les rôles et les colonnes clés des principales tables du schéma.}\normalsize

\begin{figure}[H]
    \centering
    \includegraphics[width=0.7\textwidth]{photos/ERD.png}
    \caption{Modèle Conceptuel de Données (MCD)}
    \label{fig:erd_diagram}
\end{figure}

\vspace{0.2cm}
\noindent\small\textit{Ce schéma illustre la structure de la base de données relationnelle, mettant en évidence les liaisons entre les utilisateurs, leurs équipes virtuelles, les ligues et les statistiques des joueurs réels.}
\normalsize

Le Modèle Conceptuel de Données (Figure \ref{fig:erd_diagram}) structure les relations clés suivantes :
\begin{itemize}
    \item \textbf{Utilisateurs et Équipes} : Un utilisateur (\texttt{USERS}) peut posséder plusieurs équipes (\texttt{TEAMS}), à raison d'une seule équipe par ligue.
    \item \textbf{Composition d'équipe} : La table associative \texttt{TEAM\_PLAYERS} lie les équipes virtuelles aux joueurs réels du catalogue (\texttt{PLAYERS}), représentant l'effectif sélectionné de 15 joueurs.
    \item \textbf{Performance des joueurs} : Les statistiques réelles de chaque joueur (buts, cartons, minutes jouées) et ses points calculés pour chaque journée de la compétition sont historisés dans la table \texttt{PLAYER\_SCORES}.
    \item \textbf{Batailles Head-to-Head} : Les affrontements directs entre managers se font via la table \texttt{BATTLES}, qui associe deux équipes (hôte et invitée) pour une journée (\texttt{gameweek}) précise.
\end{itemize}

\section{Modélisation structurelle --- Diagramme de classes}

Le diagramme de classes présente la structure logique du système en décrivant les classes, leurs attributs, leurs méthodes ainsi que les associations existant entre elles. Il est structuré en quatre grands packages fonctionnels :

\begin{itemize}
    \item \textbf{Comptes utilisateurs (User Accounts)} : gère les données d'authentification et de profil des utilisateurs.
    \item \textbf{Ligues, équipes \& joueurs (Leagues, Teams \& Players)} : représente le cœur du jeu de fantasy (gestion des ligues privées/publiques, composition des équipes virtuelles de 15 joueurs et catalogue des joueurs réels).
    \item \textbf{Compétition (Competition)} : gère le calendrier des journées de jeu (\textit{Gameweeks}), les scores réels des joueurs, les duels Head-to-Head (\textit{Battles}) et les prédictions communautaires.
    \item \textbf{Communication} : gère la messagerie instantanée Socket.io au sein de chaque ligue.
\end{itemize}

\begin{figure}[H]
    \centering
    \includegraphics[width=0.95\textwidth]{photos/diagramme_de_class.png}
    \caption{Diagramme de classes du système WC Fantasy}
    \label{fig:class_diagram}
\end{figure}

\vspace{0.2cm}
\noindent\small\textit{Ce diagramme de classes illustre l'organisation orientée objet du backend de WC Fantasy, montrant comment les entités collaborent pour implémenter les règles du jeu (transferts, calcul de points, duels, prédictions et messagerie).}
\normalsize

\vspace{0.5cm}
Les principales classes et leurs relations sont décrites ci-dessous :
\begin{itemize}
    \item \textbf{User} : Représente le compte d'un utilisateur. Il possède des méthodes de gestion de compte (\texttt{register()}, \texttt{login()}) et est associé à ses équipes (\texttt{Team}) ainsi qu'à ses prédictions.
    \item \textbf{Team} : Représente l'équipe virtuelle d'un utilisateur pour une ligue spécifique. Elle a un budget limité (100M) et contient exactement 15 joueurs (\texttt{Player}). Ses méthodes clés sont \texttt{addPlayer()} et \texttt{transfer()}.
    \item \textbf{League} : Représente un groupe de compétition (publique ou privée) créé par un utilisateur. Une ligue organise des batailles hebdomadaires (\texttt{Battle}) entre les équipes de ses membres.
    \item \textbf{Player} : Représente un joueur de football réel participant à la Coupe du Monde 2026.
    \item \textbf{PlayerScore} : Contient les statistiques et les points d'un joueur réel pour une journée (\texttt{Gameweek}) donnée. La méthode \texttt{calcPoints()} calcule le score selon un barème précis.
    \item \textbf{Battle} : Représente un match direct (Head-to-Head) opposant deux équipes de la même ligue pendant une journée. La méthode \texttt{resolve()} compare les scores pour désigner le vainqueur.
    \item \textbf{Prediction} : Permet aux utilisateurs de parier sur les résultats ou le meilleur buteur d'une journée afin de remporter des points bonus XP.
    \item \textbf{ChatMessage} : Gère l'historique des conversations partagées dans le chat d'une ligue.
\end{itemize}

\section{Modélisation dynamique --- Diagrammes d'activité}

Quatre diagrammes d'activité illustrent les scénarios les plus importants du système :

\subsection{Diagramme 1 : Authentification utilisateur}

Ce diagramme décrit le flux d'authentification, depuis l'ouverture de l'application jusqu'au chargement de l'écran principal. Il couvre les cas d'inscription (POST /api/auth/register) et de connexion (POST /api/auth/login), avec la vérification du token JWT en cache et la gestion des erreurs de validation.

\begin{figure}[H]
    \centering
    \includegraphics[width=0.7\textwidth]{photos/image3.png}
    \caption{Diagramme d'activité --- Authentification utilisateur}
    \label{fig:auth_diagram}
\end{figure}

\vspace{0.2cm}
\noindent\small\textit{Ce logigramme détaille les étapes séquentiellement suivies par le système lors de la tentative de connexion ou d'inscription d'un utilisateur, incluant la vérification et la génération des tokens JWT.}
\normalsize

\subsection{Diagramme 2 : Composition de l'équipe (Draft)}

Ce diagramme détaille le processus de composition d'équipe. L'utilisateur consulte le catalogue filtrable des joueurs, sélectionne un joueur, et décide de l'ajouter ou le retirer selon la disponibilité du budget et les contraintes de positions. Chaque opération est vérifiée côté serveur avec un verrou transactionnel FOR UPDATE pour éviter les race conditions.

\begin{figure}[H]
    \centering
    \includegraphics[width=0.65\textwidth]{photos/image4.png}
    \caption{Diagramme d'activité --- Composition de l'équipe (Draft)}
    \label{fig:draft_diagram}
\end{figure}

\vspace{0.2cm}
\noindent\small\textit{Ce processus décrit la logique métier exécutée lorsqu'un utilisateur modifie son équipe. Il montre comment le système vérifie le budget disponible et applique les verrous transactionnels pour éviter les erreurs concurrentes.}
\normalsize

\subsection{Diagramme 3 : Prédictions communautaires}

Avant chaque journée de jeu, les utilisateurs peuvent soumettre trois prédictions : le meilleur buteur, la nation la plus performante et le nombre total de buts. Ce diagramme montre le flux de soumission, la vérification du statut de la journée et la confirmation. Une fois la journée terminée, l'admin résout les prédictions et distribue les bonus XP.

\begin{figure}[H]
    \centering
    \includegraphics[width=0.5\textwidth]{photos/activity_predictions_v2.jpg}
    \caption{Diagramme d'activité --- Prédictions communautaires}
    \label{fig:predictions_diagram}
\end{figure}

\vspace{0.2cm}
\noindent\small\textit{Ce diagramme illustre le parcours de l'utilisateur pour soumettre une prédiction sur le vainqueur ou le meilleur buteur, ainsi que le processus de validation par l'administrateur en fin de match.}
\normalsize

\subsection{Diagramme 4 : Résolution des batailles}

En fin de journée, les scores des équipes home et away sont calculés (buts$\times$4, passes$\times$3, clean sheets$\times$4, etc.), comparés, et le résultat est enregistré dans la table battles. Une notification push est envoyée aux deux équipes, et l'historique est consultable immédiatement via GET /api/battles.

\begin{figure}[H]
    \centering
    \includegraphics[width=0.7\textwidth]{photos/image6.png}
    \caption{Diagramme d'activité --- Résolution des batailles}
    \label{fig:battles_diagram}
\end{figure}

\vspace{0.2cm}
\noindent\small\textit{Cette figure résume l'algorithme de calcul qui compare les points de deux équipes s'affrontant dans une bataille Head-to-Head, avant d'enregistrer le résultat final et de notifier les vainqueurs.}
\normalsize

\subsection{Diagramme de Séquence : Synchronisation des Scores}

Ce diagramme illustre le flux technique de mise à jour automatique des scores et points via la tâche Cron et les WebSockets :

\begin{figure}[H]
    \centering
    \includegraphics[width=0.85\textwidth]{photos/Synchronisation_Scores.png}
    \caption{Diagramme de séquence --- Processus de synchronisation des scores}
    \label{fig:sequence_scores}
\end{figure}

\vspace{0.2cm}
\noindent\small\textit{Ce diagramme de séquence décrit les échanges d'informations chronologiques entre le serveur, l'API externe (football-data) et l'application cliente lors de la mise à jour automatique et quotidienne des scores.}
\normalsize

Le processus de synchronisation (Figure \ref{fig:sequence_scores}) se déroule comme suit :
\begin{enumerate}
    \item \textbf{Déclenchement temporel} : Une tâche Cron programmée sur le serveur Node.js s'exécute quotidiennement à 03:00 UTC.
    \item \textbf{Récupération des résultats réels} : Le serveur interroge l'API externe \texttt{football-data.org} pour obtenir la liste des matchs finalisés (\texttt{FINISHED}).
    \item \textbf{Extraction et Calcul} : Pour chaque match, les statistiques individuelles de performance des joueurs sont extraites et envoyées au module \texttt{pointsCalculator} qui applique le barème métier selon la position du joueur.
    \item \textbf{Persistance} : Les points calculés sont mis à jour ou insérés dans la table \texttt{player\_scores} de la base de données PostgreSQL.
    \item \textbf{Diffusion temps réel} : Les points cumulés des équipes virtuelles et les résultats des batailles sont recalculés. Le serveur émet immédiatement un événement de notification \texttt{scores:updated} via le serveur \texttt{Socket.io} pour actualiser en temps réel l'interface client de tous les utilisateurs connectés.
\end{enumerate}

\section{Choix technologiques}

Le tableau ci-dessous récapitule chaque technologie retenue et sa justification :

\vspace{0.3cm}

\begin{table}[!htbp]
\centering
    \begin{tabularx}{\linewidth}{|l|X|}
        \hline
        \rowcolor{emsigreen}\textcolor{white}{\textbf{Technologie}} & \textcolor{white}{\textbf{Justification du choix}} \\
        \hline
        Node.js + Express.js \cite{nodejs} & Ecosystème JS unifié, asynchrone natif (I/O non bloquant), nombreux packages (jwt, bcrypt, pg, socket.io), performance élevée pour des API REST. \\
        \hline
        PostgreSQL \cite{postgresql} & SGBDR ACID complet, support JSON/JSONB pour stocker les compositions d'équipes, UUID natifs, contraintes CHECK pour l'intégrité (budget $\geq$ 0). \\
        \hline
        Redis \cite{redis} & Cache en mémoire ultra-rapide ($<$ 1ms), réduction drastique des requêtes DB répétitives, TTL automatique, pub/sub possible. \\
        \hline
        Socket.io \cite{socketio} & Abstraction WebSocket avec fallback HTTP long-polling, rooms nommées, authentification middleware JWT, reconnexion automatique. \\
        \hline
        Flutter + Dart \cite{flutter} & Un seul codebase pour iOS et Android, performances natives (moteur Skia), Riverpod pour la gestion d'état réactive, widgets riches. \\
        \hline
        Riverpod \cite{riverpod} & Gestion d'état compilée, thread-safe, testable sans BuildContext, supporte l'async avec AsyncNotifier. \\
        \hline
        JWT \cite{jwt} & Stateless, scalable horizontalement, rotation sécurisée du refresh token, révocation possible. \\
        \hline
        Brevo SMTP \cite{brevo} & Service email gratuit (300 emails/jour), fiable, supporte les templates HTML, simple à intégrer via Nodemailer \cite{nodemailer}. \\
        \hline
        Google Gemini API \cite{gemini} & LLM performant, API gratuite (niveau standard), réponses contextuelles pour le chatbot football. \\
        \hline
        football-data.org \cite{football_data} & API officielle FIFA/UEFA, couvre la Coupe du Monde, données structurées (matchs, joueurs, scores), plan gratuit disponible. \\
        \hline
    \end{tabularx}
\caption{Justification des choix technologiques}
\end{table}

\vspace{0.2cm}
\noindent\small\textit{Ce tableau expose la justification technique et fonctionnelle associée au choix de chaque brique technologique retenue pour le projet.}\normalsize

\newpage
\section*{Conclusion}

Ce chapitre a formalisé l'ensemble des besoins fonctionnels et non fonctionnels de WC Fantasy, identifié les acteurs, présenté l'architecture globale et les modèles UML qui guident le développement. Les choix technologiques retenus --- Node.js, PostgreSQL, Redis, Socket.io, Flutter --- constituent une stack moderne, cohérente et adaptée aux contraintes du projet. Le chapitre suivant présente la réalisation concrète de ce système et les résultats obtenus.
